\appendix
\section{Supporting Appendix}\label{app:support}

US exports use HMRC uktradeinfo OTS/RTS data for 2024, mapped from SITC to SIC 2007 divisions and divided by the latest non-suppressed ONS Annual Business Survey turnover. Of a \pounds 55,600 million customs-basis total, \pounds 52,500 million maps to divisions after documented exclusions, chiefly non-monetary gold and precious metals. The measured window uses May 2025--February 2026 against the mean of the same calendar months one and two years earlier. Its \RealisedAggregateFall\ per cent aggregate fall is descriptive, not causal.

\section{Shock mechanics}\label{app:mechanics}

Each employee in exposed division $j$ carries baseline displacement risk $s_j$. The primary estimator realises that risk through the balanced systematic candidate draw of Section~\ref{sec:method-margins}, which softly matches expected industry wage-bill and headcount targets; the Bernoulli comparator instead selects each employee independently with probability $s_j$, so its weighted displaced counts and earnings meet targets in expectation rather than exactly. Selected workers have employment income, hours, pension contributions and statutory pay zeroed. Wage cuts remove $s_j$ of each exposed worker's earnings deterministically. Reallocation uses the same selected workers, assigns a services destination, and applies the calibrated earnings penalty.

UC claiming is redrawn after every margin for benefit units whose circumstances changed and that are entitled post-shock. The LCWRA override sets the benefit-unit element to the maximum of its stored baseline value and one statutory annual element, preventing duplicate payment.

\section{Monte Carlo and measures}\label{app:mc}

Central direct and measured scenarios use \ProductionDraws\ assignments. Every displayed $\pm$ value is an assignment SD, not a standard error or confidence interval. The machine-readable artifacts additionally store numerical Monte Carlo standard error as assignment SD divided by $\sqrt{n}$; it measures precision of the simulated mean only. Survey-design, tariff-response and model-parameter uncertainty are not estimated. Five-assignment mechanism, demographic, reallocation, duration and sensitivity-grid exercises, and the ten-assignment supply-chain exercise, are exploratory and are excluded from the headline evidence; the affected-household poverty exercise is scored at \WelfarePovDraws\ assignments and the take-up grid at \TakeupSeeds\ assignments, and both are reported in the main text. One further exception to the $\pm$ convention: the zero-hours diagnostic quotes Monte Carlo standard errors over its seeds, and says so where the values appear.
The mixed-adjustment scenario surface likewise uses five common assignments
per cell. Its machine-readable cells and draws are in
\path{results/scenario_testing.json} and \path{results/scenario_testing.csv}.

Reported relative BHC poverty changes hold the poverty line frozen at 60 per cent of the baseline equivalised median; absolute BHC poverty likewise uses 60 per cent of the baseline equivalised median as a fixed line. The cushioning rate is $1-\Delta Y/\Delta E$, where $\Delta E$ is gross employment-earnings loss and $\Delta Y$ is household disposable-income loss.

\paragraph{Thin-support cell.} The OBR three-month duration-equivalent cell is omitted from the main bounds table because its displacement draw rests on \OBRThreeMonthEffectiveLossRecords\ effective record contributions, with the largest record supplying \OBRThreeMonthMaxRecordLossShare\ per cent of the loss on average. For completeness:

\begin{center}
\small
\begin{tabular}{llrrrrr}
\toprule
Anchor & Months & Margin & Gross loss & Exchequer & Cushioning & Poverty \\
& & & (\pounds m) & (\pounds m) & (\%) & (pp) \\
\midrule
\SubmissionScenarioRowsThin
\bottomrule
\end{tabular}
\end{center}

These rows are sensitivity evidence only, not population estimates.

\paragraph{Reconciling cushioning and the Exchequer effect.} The Exchequer effect is the change in the aggregate simulated government balance, not the household-side transfer implied by the cushioning identity. The identity implies a household-side offset of $(1-\Delta Y/\Delta E)\times\Delta E=\Delta E-\Delta Y$: \pounds\FullWageImpliedOffset\ million for the full-schedule wage cut and \pounds\FullDisplacedImpliedOffset\ million for displacement (computed as $\Delta E-\Delta Y$ within each draw and then averaged, not as mean cushioning times mean gross loss), against the \ProductionDraws-assignment Bernoulli run's Exchequer costs of \pounds\FullWageExchequer\ million and \pounds\FullDisplacedExchequer\ million. The reconciliation is internally coherent on that vintage; Table~\ref{tab:cushioning} reports the same displacement scenario under the \SubmissionDraws-assignment balanced design, at \pounds\BalancedUnitExchequer\ million. The wedge is revenue lost outside employee disposable income: employer National Insurance contributions fall with the employee wage bill under every margin, and for displaced workers the zeroing of employee pension contributions and salary-sacrifice arrangements changes taxable pay relative to the gross-earnings accounting in $\Delta E$. The government balance also absorbs any other simulated fiscal flows that respond to the shock. The household-side offset therefore understates the Exchequer cost by \pounds\FullDisplacedExchequerWedge--\FullWageExchequerWedge\ million in the central full-schedule scenarios, and the two statistics answer different questions: the cushioning rate measures household income protection per pound of employee earnings lost, while the Exchequer effect measures the total first-round cost to the public finances.

\section{Factorial adjustment scenarios}\label{app:factorial}

Figure~\ref{fig:scenario-testing} separates the aggregate export-demand
calibration from the assumed composition of labour-market adjustment. Each
cell uses the full tariff schedule and five common assignments. The wage-cut
endpoint is stable: cushioning is
\ScenarioWageEndpointMin--\ScenarioWageEndpointMax\ per cent over the four
calibrations. Moving towards displacement generally lowers cushioning, but
the small-shock cells are especially noisy because a few high-weight FRS
records dominate five Bernoulli assignments. These cells are stress
scenarios, not estimates or confidence regions. The qualitative result is
that fiscal cost scales mainly with the imposed gross shock, while the
adjustment mix changes the share absorbed by taxes and transfers.

\section{Duration scope}\label{app:durationtakeup}

The submission design reports three-, six- and twelve-month
\emph{duration-equivalent annual stresses}. It scales the probability of a
coherent full-year displaced state rather than combining partial-year earnings
with full-year unemployment status. These rows preserve expected worker-month
loss but are not unemployment-spell simulations. Actual within-year tax, UC,
waiting-period and re-employment timing requires a monthly model and is not
claimed here.

\section{Monthly versus annual UC accounting}\label{app:monthlyuc}

The annual model represents an $m$-month unemployment spell as a
probability $m/12$ of a coherent full-year displaced state. This appendix
bounds the resulting bias with statutory 2026 parameters for a
representative single displaced worker aged 25 or over, without children,
housing element or capital above the lower limit, at the exposed-sector FRS
mean wage of \pounds 48{,}272. UC entitlement per worker-month is identical
by construction: $m$ months at the full standard allowance
(\pounds\MonthlyUCStandardAllowance\ a month; earnings above the taper-out
point in the working months) equals probability $m/12$ of twelve months.
The differences run through income tax and timing. Because the tax schedule
is convex, zeroing a full year of earnings forgoes relief at the worker's
top marginal slices, whereas a partial-year earner sheds basic- and
higher-rate slices: at this wage, the correct partial-year PAYE relief for
a three-month spell is \pounds\MonthlyUCTaxMonthlyThree\ against
\pounds\MonthlyUCTaxAnnualThree\ in the annual equivalent, so the annual
model \emph{understates} total cushioning for a three- or six-month spell
by about \MonthlyUCCushionGap\ percentage points
(\MonthlyUCCushionMonthly\ versus \MonthlyUCCushionAnnual\ per cent) and is
exact at twelve months. The five-week UC payment wait shifts timing, not
annual entitlement. The primary 12-month contrast in
Table~\ref{tab:cushioning} is therefore unaffected, and the shorter
duration-equivalent rows are conservative for the displacement margin:
correcting the annualisation would raise displacement cushioning at those
scales and narrow, not widen, the reported gap. Parameters and the full
accounting are stored in \path{results/referee_fixes.json}.
